% Generated draft: requires mathematical and source review.
\documentclass[11pt]{article}
\usepackage[margin=25mm]{geometry}
\usepackage{fontspec}
\setmainfont{Latin Modern Roman}
\usepackage{amsmath,amssymb,mathtools,unicode-math}
\setmathfont{Latin Modern Math}
\usepackage{enumitem,needspace,xurl,hyperref}
\setlength{\parindent}{0pt}
\setlength{\parskip}{5pt}
\setlength{\emergencystretch}{3em}
\newcommand{\N}{\mathbb{N}}
\newcommand{\Z}{\mathbb{Z}}
\newcommand{\Q}{\mathbb{Q}}
\newcommand{\R}{\mathbb{R}}
\newcommand{\C}{\mathbb{C}}
\newcommand{\F}{\mathbb{F}}
\newcommand{\A}{\mathbb{A}}
\newcommand{\PP}{\mathbb{P}}
\newcommand{\E}{\mathbb{E}}
\newcommand{\eps}{\varepsilon}
\newcommand{\dd}{\,\mathrm{d}}
\begin{document}
\section*{UnsolvedMath ID: 30005220}
\textbf{Source problem code:} catalogue entry
\subsection{Definitions and mathematical statement}
Let \(p\) be a prime, let \(G\) be a finite group, and let \(P\leq G\) be a Sylow \(p\)-subgroup, so that \(|P|\) is the largest power of \(p\) dividing \(|G|\). Write \(\operatorname{Irr}(G)\) for the irreducible complex characters of \(G\). A character \(\chi\in\operatorname{Irr}(G)\) has \(p'\)-degree if \(p\nmid\chi(1)\).

For a positive integer \(n\), put \(\zeta_n=e^{2\pi i/n}\) and \(\Q_n=\Q(\zeta_n)\). Define
\[
\Q(\chi)=\Q\bigl(\chi(g):g\in G\bigr),\qquad
\chi_P=\operatorname{Res}^{G}_{P}\chi,\qquad
\Q(\chi_P)=\Q\bigl(\chi(x):x\in P\bigr).
\]
Here \(\chi_P\) need not be irreducible. Interpret the conductor of \(\chi\) in the standard cyclotomic sense as
\[
c(\chi)=\min\{n\geq1:\Q(\chi)\subseteq\Q_n\}.
\]
Write \(c(\chi)=p^a m\), where \(a\in\Z_{\geq0}\), \(m\geq1\), and \(p\nmid m\). With this convention, the restriction field lies in \(\Q_{p^a}\), so the following field degree is defined; in particular, \(\Q_{p^0}=\Q_1=\Q\).

\textbf{Question from the supplied corrected statement.} Is it true that, for every prime \(p\), every finite group \(G\), every Sylow \(p\)-subgroup \(P\) of \(G\), and every \(\chi\in\operatorname{Irr}(G)\) of \(p'\)-degree with conductor \(p^a m\) as above,
\[
p\nmid[\Q_{p^a}:\Q(\chi_P)]\,?
\]

\textbf{Source qualifications.} The conductor convention above makes explicit terminology left undefined in the supplied statement. The general-prime quantifier follows the supplied corrected statement. The supplied original excerpt is corrupted and ends with an instruction to fix \(p=2\); it does not contain the displayed question in recoverable form. Thus the intended scope in the original report cannot be confirmed from this excerpt alone.

The surrounding excerpt motivates the question through characters of \(p'\)-degree and the McKay conjecture. It also states that character fields are abelian extensions of \(\Q\), and that, for \(\Q\subseteq F\subseteq\Q_n\), the semidirect product \(C_n\rtimes\operatorname{Gal}(\Q_n/F)\), using the cyclotomic action on the cyclic group \(C_n\), has an irreducible character with field of values \(F\). These are supplied contextual claims, not additional hypotheses of the question.
\subsection{Short English statement}
For an irreducible character whose degree is prime to \(p\), take the \(p\)-part of its cyclotomic conductor. Must the field generated by its values on a Sylow \(p\)-subgroup have index prime to \(p\) in the corresponding cyclotomic field? This draft follows the supplied general-prime formulation; the corrupted original excerpt also mentions specializing to \(p=2\).
\subsection{Sources}
\begin{itemize}
\item Supplied local UnsolvedMath registry record, permanent ID \(30005220\), source code \texttt{catalogue entry}, title \emph{Sylow Restrictions and Character Fields of Values}, in the Oberwolfach Reports Open Problems collection. The project registry paths are \texttt{data/unsolvedmath/problems.json} and \texttt{catalog/problems.json}. No individual proposer is supplied; the recorded proposed year is 2022.
\item Supplied report citation: \emph{Character Theory and Categorification (2023)}, Oberwolfach Reports, pp.\ 2267--2333, DOI \texttt{10.4171/owr/2022/39}. \url{https://doi.org/10.4171/owr/2022/39}. The supplied citation uses 2023 while the DOI and report designation refer to 2022; these details are retained as supplied.
\item Supplied literature source: Fayers--Kleshchev--Morotti, spin RoCK blocks. \url{https://doi.org/10.1112/jlms.12852}. Its precise relevance to the displayed question is not established by the supplied record.
\item Supplied literature source: Oberwolfach Report 39/2022. \url{https://ems.press/content/serial-article-files/46974}.
\item Literature assessment dated \texttt{2026-08-22}: the supplied curation labels the problem open and reports motivations and partial results, but no general theorem resolving the relation between character restrictions and fields of values was found. This is supplied dated curation, not an independently verified literature assessment. The record itself requests independent verification of time-sensitive status and cited claims.
\item The supplied statement status is \texttt{corrected\_verified}; its verification note says that the official report was checked and the question isolated with repaired notation. This records the curator's assertion, not an independent source check performed for this draft.
\end{itemize}
\end{document}
